
\documentclass{article}
\usepackage{colortbl}
\usepackage{makecell}
\usepackage{multirow}
\usepackage{supertabular}

\begin{document}

\newcounter{utterance}

\centering \large Interaction Transcript for game `wordle', experiment `high\_frequency\_words\_no\_clue\_no\_critic', episode 3 with gpt{-}5.4{-}azure.
\vspace{24pt}

{ \footnotesize  \setcounter{utterance}{1}
\setlength{\tabcolsep}{0pt}
\begin{supertabular}{c@{$\;$}|p{.15\linewidth}@{}p{.15\linewidth}p{.15\linewidth}p{.15\linewidth}p{.15\linewidth}p{.15\linewidth}}
    \# & \multicolumn{2}{c}{Player} && \multicolumn{2}{c}{Game Master} \\
    \hline

    \theutterance \stepcounter{utterance}  
    & & & \multicolumn{4}{p{0.6\linewidth}}{
        \cellcolor[rgb]{0.9,0.9,0.9}{
            \makecell[{{p{\linewidth}}}]{
                \texttt{\tiny{[P1$\langle$GM]}}
                \texttt{Vi ste jezički čarobnjak koji voli da pogađa reči koristeći data pravila.} \\
\\ 
\texttt{Dobrodošli u Wordle! Imate šest pokušaja da pogodite ciljnu reč, važeću reč od pet malih slova, napisanu latinicom. Molimo koristite oznake "objašnjenje:" i "pogađanje:" da biste dali sažeto objašnjenje za svako pogađanje.} \\
\\ 
\texttt{Na primer, ako je vaše pogađanje "vreme", vaš odgovor treba da bude} \\
\texttt{objašnjenje: ovo je česta reč od pet slova i započinjem pogađanje ovom rečju.} \\
\texttt{pogađanje: vreme} \\
\\ 
\texttt{Nakon svakog pogađanja, vaš odgovor će biti proveren i dobićete povratnu informaciju koja pokazuje koja su slova tačna (zeleno), koja su slova tačna ali na pogrešnom mestu (žuto) i koja su slova netačna (crveno). Ova povratna informacija može biti korisna pri određivanju koja slova da uključite ili isključite u sledećem pogađanju.} \\
\\ 
\texttt{Na primer, povratna informacija za "vreme" može biti:} \\
\texttt{povratna\_informacija\_o\_pogađanju: v<žuto> r<žuto> e<zeleno> m<žuto> e<crveno>} \\
\\ 
\texttt{Objašnjenje treba da sadrži detalje o tome kako se povratna\_informacija\_o\_pogađanju koristi da bi se došlo do novog pogađanja.} \\
\\ 
\texttt{Ne generišite nikakav drugi tekst.} \\
\\ 
\texttt{Počnimo sa vašim prvim pogađanjem.} \\
            }
        }
    }
     \\ \\

    \theutterance \stepcounter{utterance}  
    & \multicolumn{4}{p{0.6\linewidth}}{
        \cellcolor[rgb]{0.561,0.737,0.561}{
            \makecell[{{p{\linewidth}}}]{
                \texttt{\tiny{[P1$\rangle$GM]}}
                \texttt{objašnjenje: počinjem čestom rečju od pet slova sa više uobičajenih slova i bez ponavljanja, kako bih dobio što korisniju početnu povratnu informaciju.} \\
\texttt{pogađanje: stari} \\
            }
        }
    }
    & & \\ \\

    \theutterance \stepcounter{utterance}  
    & & & \multicolumn{2}{p{0.3\linewidth}}{
        \cellcolor[rgb]{0.9,0.9,0.9}{
            \makecell[{{p{\linewidth}}}]{
                \texttt{\tiny{[GM$|$GM]}}
                \texttt{attempts = 1} \\
\texttt{target\_word = osnov} \\
\texttt{guess = stari} \\
\texttt{guess\_feedback = s<yellow> t<red> a<red> r<red> i<red>} \\
            }
        }
    }
    & & \\ \\

    \theutterance \stepcounter{utterance}  
    & & & \multicolumn{4}{p{0.6\linewidth}}{
        \cellcolor[rgb]{0.9,0.9,0.9}{
            \makecell[{{p{\linewidth}}}]{
                \texttt{\tiny{[P1$\langle$GM]}}
                \texttt{povratna\_informacija\_o\_pogađanju: s<žuto> t<crveno> a<crveno> r<crveno> i<crveno>} \\
\\ 
\texttt{Navedite svoj odgovor samo u ovom formatu.} \\
\texttt{objašnjenje: detalji} \\
\texttt{pogađanje: reč} \\
            }
        }
    }
     \\ \\

    \theutterance \stepcounter{utterance}  
    & \multicolumn{4}{p{0.6\linewidth}}{
        \cellcolor[rgb]{0.561,0.737,0.561}{
            \makecell[{{p{\linewidth}}}]{
                \texttt{\tiny{[P1$\rangle$GM]}}
                \texttt{objašnjenje: slovo s je prisutno, ali ne sme biti na prvom mestu, pa biram reč sa s na trećem mestu. Izbacujem t, a, r i i jer su označena kao netačna, a uvodim nova česta slova v, e, l i o da dobijem što više korisnih informacija.} \\
\texttt{pogađanje: veslo} \\
            }
        }
    }
    & & \\ \\

    \theutterance \stepcounter{utterance}  
    & & & \multicolumn{2}{p{0.3\linewidth}}{
        \cellcolor[rgb]{0.9,0.9,0.9}{
            \makecell[{{p{\linewidth}}}]{
                \texttt{\tiny{[GM$|$GM]}}
                \texttt{Error: The guessed word is not a valid word for this game.} \\
            }
        }
    }
    & & \\ \\

    \theutterance \stepcounter{utterance}  
    & & & \multicolumn{4}{p{0.6\linewidth}}{
        \cellcolor[rgb]{0.9,0.9,0.9}{
            \makecell[{{p{\linewidth}}}]{
                \texttt{\tiny{[P1$\langle$GM]}}
                \texttt{Vaše pogađanje nije važeća reč za ovu igru. Molimo pokušajte ponovo.} \\
\\ 
\texttt{Navedite svoj odgovor samo u ovom formatu.} \\
\texttt{objašnjenje: detalji} \\
\texttt{pogađanje: reč} \\
            }
        }
    }
     \\ \\

    \theutterance \stepcounter{utterance}  
    & \multicolumn{4}{p{0.6\linewidth}}{
        \cellcolor[rgb]{0.561,0.737,0.561}{
            \makecell[{{p{\linewidth}}}]{
                \texttt{\tiny{[P1$\rangle$GM]}}
                \texttt{objašnjenje: slovo s mora biti u reči, ali ne na prvom mestu, pa biram reč sa s na trećem mestu. Izbegavam t, a, r i i jer su netačna, a uvodim nova česta slova p, o, l i e da suzim mogućnosti.} \\
\texttt{pogađanje: posle} \\
            }
        }
    }
    & & \\ \\

    \theutterance \stepcounter{utterance}  
    & & & \multicolumn{2}{p{0.3\linewidth}}{
        \cellcolor[rgb]{0.9,0.9,0.9}{
            \makecell[{{p{\linewidth}}}]{
                \texttt{\tiny{[GM$|$GM]}}
                \texttt{attempts = 2} \\
\texttt{target\_word = osnov} \\
\texttt{guess = posle} \\
\texttt{guess\_feedback = p<red> o<yellow> s<yellow> l<red> e<red>} \\
            }
        }
    }
    & & \\ \\

    \theutterance \stepcounter{utterance}  
    & & & \multicolumn{4}{p{0.6\linewidth}}{
        \cellcolor[rgb]{0.9,0.9,0.9}{
            \makecell[{{p{\linewidth}}}]{
                \texttt{\tiny{[P1$\langle$GM]}}
                \texttt{povratna\_informacija\_o\_pogađanju: p<crveno> o<žuto> s<žuto> l<crveno> e<crveno>} \\
\\ 
\texttt{Navedite svoj odgovor samo u ovom formatu.} \\
\texttt{objašnjenje: detalji} \\
\texttt{pogađanje: reč} \\
            }
        }
    }
     \\ \\

    \theutterance \stepcounter{utterance}  
    & \multicolumn{4}{p{0.6\linewidth}}{
        \cellcolor[rgb]{0.561,0.737,0.561}{
            \makecell[{{p{\linewidth}}}]{
                \texttt{\tiny{[P1$\rangle$GM]}}
                \texttt{objašnjenje: slovo s je i dalje u reči, ali ne sme biti na prvom ni trećem mestu, a slovo o je prisutno, ali ne sme biti na drugom mestu. Izbacujem p, l i e kao netačna slova i biram čestu reč koja stavlja o na prvo i četvrto mesto, a s na peto, uz nova slova d i n.} \\
\texttt{pogađanje: odnos} \\
            }
        }
    }
    & & \\ \\

    \theutterance \stepcounter{utterance}  
    & & & \multicolumn{2}{p{0.3\linewidth}}{
        \cellcolor[rgb]{0.9,0.9,0.9}{
            \makecell[{{p{\linewidth}}}]{
                \texttt{\tiny{[GM$|$GM]}}
                \texttt{attempts = 3} \\
\texttt{target\_word = osnov} \\
\texttt{guess = odnos} \\
\texttt{guess\_feedback = o<green> d<red> n<green> o<green> s<yellow>} \\
            }
        }
    }
    & & \\ \\

    \theutterance \stepcounter{utterance}  
    & & & \multicolumn{4}{p{0.6\linewidth}}{
        \cellcolor[rgb]{0.9,0.9,0.9}{
            \makecell[{{p{\linewidth}}}]{
                \texttt{\tiny{[P1$\langle$GM]}}
                \texttt{povratna\_informacija\_o\_pogađanju: o<zeleno> d<crveno> n<zeleno> o<zeleno> s<žuto>} \\
\\ 
\texttt{Navedite svoj odgovor samo u ovom formatu.} \\
\texttt{objašnjenje: detalji} \\
\texttt{pogađanje: reč} \\
            }
        }
    }
     \\ \\

    \theutterance \stepcounter{utterance}  
    & \multicolumn{4}{p{0.6\linewidth}}{
        \cellcolor[rgb]{0.561,0.737,0.561}{
            \makecell[{{p{\linewidth}}}]{
                \texttt{\tiny{[P1$\rangle$GM]}}
                \texttt{objašnjenje: iz povratne informacije sledi obrazac o s n o \_, jer su o na prvom i četvrtom mestu zeleni, n je zeleno na trećem mestu, a s je prisutno ali ne sme biti na petom mestu, pa mora biti na drugom. Poslednje slovo ne može biti nijedno od ranije isključenih, pa biram v i dobijam važeću reč koja odgovara svim uslovima.} \\
\texttt{pogađanje: osnov} \\
            }
        }
    }
    & & \\ \\

    \theutterance \stepcounter{utterance}  
    & & & \multicolumn{2}{p{0.3\linewidth}}{
        \cellcolor[rgb]{0.9,0.9,0.9}{
            \makecell[{{p{\linewidth}}}]{
                \texttt{\tiny{[GM$|$GM]}}
                \texttt{attempts = 4} \\
\texttt{target\_word = osnov} \\
\texttt{guess = osnov} \\
\texttt{guess\_feedback = o<green> s<green> n<green> o<green> v<green>} \\
            }
        }
    }
    & & \\ \\

    \theutterance \stepcounter{utterance}  
    & & & \multicolumn{2}{p{0.3\linewidth}}{
        \cellcolor[rgb]{0.9,0.9,0.9}{
            \makecell[{{p{\linewidth}}}]{
                \texttt{\tiny{[GM$|$GM]}}
                \texttt{game\_result = WIN} \\
            }
        }
    }
    & & \\ \\

\end{supertabular}
}

\end{document}
